%\pagebreak
\section{List Item Expressions}
\label{sec:list_item_expressions}
\index{directive syntax!list item expressions}
\index{list item expressions}
\index{containing array}
\index{containing structure}
\index{base variable}
\index{base array}
\index{base pointer}
\index{named pointer}
\index{array base}

Directives, clauses, and modifiers may accept a list argument that consists of
one or more comma-separated list items.  An expression that forms a list item must
satisfy syntactic requirements defined by the specification. When the list item
refers to a data object in memory, these requirements are often defined using
terms that refer to various parts of the expression. These terms include:
\emph{containing array}, \emph{containing structure}, \emph{base variable},
\emph{base array}, \emph{base pointer}, \emph{named pointer}, and \emph{array
base}.

In the examples provided, identifiers starting with \ucode{p} are pointers,
identifiers starting with \ucode{x} are arrays, identifiers starting with
\ucode{s} are structures, and identifiers starting with \ucode{y} are scalars.
For generality, we use a unified description for list item expressions in C,
C++, and Fortran. To this end, we treat a Fortran array element or array
section designator as composed of an array type subexpression (the array) to
which an array subscript operator is applied that contains its subscript list,
and we treat a Fortran structure component designator as composed of a
structure type subexpression (the structure) to which a structure element
operator is applied. We define a \emph{direct-selection operator} as an array
subscript operator (\ucode{[...]} in C/C++ or \ucode{(...)} in Fortran) or a
structure element operator (\ucode{.} in C/C++ or \ucode{\%} in Fortran) that
is not applied to a pointer variable, after transforming (for C/C++)
expressions of the form \ucode{p->x} to \ucode{(*p).x}.  And we define an
\emph{indirect-selection operator} as an array subscript operator
(\ucode{[...]} in C/C++ or \ucode{(...)} in Fortran) or structure element
operator (\ucode{->} in C/C++ and \ucode{\%} in Fortran) that is applied to a
pointer variable, after transforming (for C/C++) expressions of the form
\ucode{*p} to \ucode{*(p+0)} and then expressions of the form \ucode{*(p+i)} to
\ucode{p[i]}, where \ucode{p} is a pointer and \ucode{i} is an integer type.
So, \ucode{x[i]} uses a direct-selection operator if \ucode{x} is an array and
not a pointer, while \ucode{p[i]} uses an indirect-selection operator if
\ucode{p} is a pointer.

In the following explanations, a subexpression is assumed to be an expression
that can produce a value; in particular, for structures the name of a structure
component that follows a structure element operator is not considered a
subexpression.

\subsection{Containing Arrays}

The containing array of an expression that represents a data object is any
unsubscripted subexpression of array type to which a sequence of zero or more
direct-selection operators are applied to yield the data object expression. So:

\begin{itemize}
\item the containing array of \ucode{x1} is \ucode{x1};

\item the containing array of \ucode{x1[i]} is \ucode{x1} for C/C++ and of
      \ucode{x1(i)} is \ucode{x1} for Fortran;

\item the containing array of \ucode{x1[i:4]} is \ucode{x1} for C/C++ and of
      \ucode{x1(i:i+3)} is \ucode{x1} for Fortran;

\item the containing array of \ucode{p1->x1[i].y1} is \ucode{p1->x1} for C/C++
      and of \ucode{p1\%x1(i)\%y1} is \ucode{p1\%x1} for Fortran;

\item the containing arrays of \ucode{p1->x1[i].x2} are \ucode{p1->x1} and
      \ucode{p1->x1[i].x2} for C/C++ and of \ucode{p1\%x1(i)\%x2} are
      \ucode{p1\%x1} and \ucode{p1\%x1(i)\%x2} for Fortran;

\item the containing arrays of \ucode{p1->x1[i].x2[j]} are \ucode{p1->x1} and
      \ucode{p1->x1[i].x2} for C/C++ and of \ucode{p1\%x1(i)\%x2\%(j)} are
      \ucode{p1\%x1} and \ucode{p1\%x1(i)\%x2} for Fortran;

\item the containing arrays of \ucode{p1->x1[i].x2[j:4]} are \ucode{p1->x1} and
      \ucode{p1->x1[i].x2} for C/C++ and of \ucode{p1\%x1(i)\%x2(j:j+3)} are
      \ucode{p1\%x1} and \ucode{p1\%x1(i)\%x2} for Fortran;

\item the containing array of \ucode{p1->p2[i]->x1} is \ucode{p1->p2[i]->x1}
      for C/C++ and of \ucode{p1\%p2(i)\%x1} is \ucode{p1\%p2(i)\%x1} for Fortran;

\item the containing array of \ucode{p1->p2[i]->x1[j]} is \ucode{p1->p2[i]->x1}
      for C/C++ and of \ucode{p1\%p2(i)\%x1(j)} is \ucode{p1\%p2(i)\%x1} for
      Fortran;

\item the containing array of \ucode{p1->p2[i]->x1[j:4]} is
      \ucode{p1->p2[i]->x1} for C/C++ and of \ucode{p1\%p2(i)\%x1(j:j+3)} is
      \ucode{p1\%p2(i)\%x1} for Fortran.
\end{itemize}


\subsection{Containing Structures}

The containing structure of an expression that represents a data object is any
subexpression of structure type to which a sequence of zero or more
direct-selection operators are applied to yield the data object expression.
Additionally, for Fortran, the first operator may be an indirect-selection
operator that is applied to a structure variable with the \bcode{POINTER} attribute, so
that the containing structure is treated as the pointee of the variable. So:

\begin{itemize}
\item the containing structure of \ucode{s1} is \ucode{s1};

\item the containing structure of \ucode{s1.y} is \ucode{s1} for C/C++ and of
      \ucode{s1\%y} is \ucode{s1} for Fortran;

\item the containing structures of \ucode{s1.x1[i].s2} are \ucode{s1.x1[i].s2},
      \ucode{s1.x1[i]}, and \ucode{s1} for C/C++ and of \ucode{s1\%x1(i)\%s2}
      are \ucode{s1\%x1(i)\%s2}, \ucode{s1\%x1(i)}, and \ucode{s1} for
      Fortran;

\item the containing structures of \ucode{p1->x1[i].y} are \ucode{p1->x1[i]} and
      \ucode{*p1} for C/C++ and of \ucode{p1\%x1(i)\%y} are \ucode{p1\%x1(i)} and
      \ucode{p1} for Fortran.
\end{itemize}


% <Add example illustrating one or more rules or restrictions pertaining to containing structures.>

\subsection{Base Variables}

If an expression that represents a data object or one of its containing arrays
or containing structures is denoted by a variable that is a base language
identifier, that variable is a base variable of that expression. So:

\begin{itemize}
\item the base variable of \ucode{y} is \ucode{y};

\item the base variable of \ucode{s.y} is \ucode{s} for C/C++ and of
      \ucode{s\%y} is \ucode{s} for Fortran;

\item the base variable of \ucode{x[i].s.y} is \ucode{x} for C/C++ and of
      \ucode{x(i)\%s\%y} is \ucode{x} for Fortran.
\end{itemize}

% <Add example illustrating rules or restrictions pertaining to base variables.>

\subsection{Base Arrays}

If an expression that represents a data object has one or more containing
arrays, its base array is the containing array that is contained by all of
its other containing arrays. So:

\begin{itemize}
\item the base array of \ucode{x1} is \ucode{x1};

\item the base array of \ucode{x1[i]} is \ucode{x1} for C/C++ and of
      \ucode{x1(i)} is \ucode{x1} for Fortran;

\item the base array of \ucode{x1[i:4]} is \ucode{x1} for C/C++ and of
      \ucode{x1(i:i+3)} is \ucode{x1} for Fortran;

\item the base array of \ucode{p1->x1[i].y1} is \ucode{p1->x1} for C/C++ and of
      \ucode{p1\%x1(i)\%y1} is \ucode{p1\%x1} for Fortran;

\item the base array of \ucode{p1->x1[i].x2} is \ucode{p1->x1.x2} for C/C++ and of
      \ucode{p1\%x1(i)\%x2} is \ucode{p1\%x1\%x2} for Fortran;

\item the base array of \ucode{p1->x1[i].x2[j]} is \ucode{p1->x1.x2} for C/C++ and
      of \ucode{p1\%x1(i)\%x2\%(j)} is \ucode{p1\%x1\%x2} for Fortran;

\item the base array of \ucode{p1->x1[i].x2[j:4]} is \ucode{p1->x1.x2} for C/C++
      and of \ucode{p1\%x1(i)\%x2(j:j+3)} is \ucode{p1\%x1\%x2} for Fortran;

\item the base array of \ucode{p1->p2[i]->x1} is \ucode{p1->p2[i]->x1} for
      C/C++ and of \ucode{p1\%p2(i)\%x1} is \ucode{p1\%p2(i)\%x1} for
      Fortran;

\item the base array of \ucode{p1->p2[i]->x1[j]} is \ucode{p1->p2[i]->x1} for
      C/C++ and of \ucode{p1\%p2(i)\%x1(j)} is \ucode{p1\%p2(i)\%x1} for
      Fortran;

\item the base array of \ucode{p1->p2[i]->x1[j:4]} is \ucode{p1->p2[i]->x1} for
      C/C++ and of \ucode{p1\%p2(i)\%x1(j:j+3)} is \ucode{p1\%p2(i)\%x1} for
      Fortran.
\end{itemize}

% <Add example illustrating rules or restrictions pertaining to base arrays.>

\subsection{Base Pointers}

The base pointer of an expression that represents a data object is a
subexpression that is a pointer to which an indirect-selection operator is
applied to yield the data object expression or a containing array or containing
structure of the data object expression. Additionally, for Fortran if the data
object is expressed as an associated data pointer, then the data pointer is the
base pointer of the data object. So:

\begin{itemize}
\item the base pointer of \ucode{*p} (equivalently \ucode{p[0]}) is \ucode{p}
      for C/C+ and of \ucode{p} is \ucode{p} for Fortran;

\item the base pointer of \ucode{*(p+i)} (equivalently \ucode{p[i]}) is
      \ucode{p} for C/C++;

\item the base pointer of \ucode{p[i]} is \ucode{p} for C/C++ and of
      \ucode{p(i)} is \ucode{p} for Fortran;

\item the base pointer of \ucode{p[i:4]} is \ucode{p} for C/C++ and of
      \ucode{p(i:i+3)} is \ucode{p} for Fortran;

\item the base pointer of \ucode{p1->p2->s1.y} is \ucode{p1->p2} for C/C++ and
      of \ucode{p1\%p2\%s1\%y} is \ucode{p1\%p2} for Fortran.
\end{itemize}

% <Add example illustrating rules or restrictions pertaining to base pointers.>

\subsection{Named Pointers}

If an expression that represents a data object has a base pointer then its
named pointers include that base pointer and any base pointers of its named
pointers. So:

\begin{itemize}
\item the named pointers of \ucode{p1[i].p2->y} are \ucode{p1[i].p2} and
      \ucode{p1} for C/C++ and of \ucode{p1(i)\%p2\%y} are \ucode{p1(i)\%p2}
      and \ucode{p1} for Fortran;

\item the named pointers of \ucode{\ucode{p}1[i][j]} for C/C++, where
      \ucode{p1} is a pointer to an array pointers, is \ucode{p1[i]} and
      \ucode{p1}.
\end{itemize}

% <Add example illustrating rules or restrictions pertaining to named pointers.>

\subsection{Array Bases}

The array base of an array section or array element is its base array if it has
one, and otherwise it is its base pointer if it has one. So:

\begin{itemize}
\item the array base of \ucode{s1.x1[i].x2[j]} is \ucode{s1.x1} for C/C++ and 
      of \ucode{s1\%x1(i)\%x2(j)} is \ucode{s1\%x1} for Fortran;

\item the array base of \ucode{s1.x1[i].p1[j]} is \ucode{s1.x1[i].p1} for C/C++ and 
      of \ucode{s1\%x1(i)\%p2(j)} is \ucode{s1\%x1(i)\%p1} for Fortran.
\end{itemize}

% <Add example illustrating rules or restrictions pertaining to array bases.>
